\documentclass{article}
\usepackage[T1]{fontenc}
\usepackage{lmodern}
\usepackage{graphicx}
\usepackage{amsmath,amssymb}
\usepackage[a4paper,margin=2cm]{geometry}
\usepackage[absolute,overlay]{textpos}
\setlength{\TPHorizModule}{1mm}
\setlength{\TPVertModule}{1mm}
\usepackage[indonesian]{babel}
\usepackage{hyperref}
\setlength{\emergencystretch}{2em}
% unit: Halaman 51

\begin{document}

% === Halaman 51 (python/native) ===

Aplikasi ECC

• Banyak piranti yang berukuran kecil dan memiliki keterbatasan 

memori dan kemampuan pemrosesan.

• Di mana kita dapat menerapkan ECC?

• Piranti komunikasi nirkabel

• Smart cards

• Web server yang membutuhkan penangangan  banyak sesi enkripsi

• Sembarang aplikasi yang membutuhkan keamanan  tetapi memiliki 

kekurangan dalam power, storage and kemampuan komputasi adalah 

potensial memerlukan ECC

*) Sumber bahan: Debdeep Mukhopadhyay, Elliptic Curve Cryptography ,

 Dept of Computer Sc and Engg IIT Madras

51

\end{document}
